% ===================================================================
%  TABELO N-RO 10 — La maro. — La haveno.
%  Traduction 2026 d'apres texto/io, controlee sur texto/fr et
%  texto/es. Consigne : voir texto/eo/10-tabelo-01.tex.
%
%  Le francais decale sa numerotation d'un cran a partir de l'alinea 8.
%  C'est une coquille de l'imprimeur francais, conservee dans la
%  transcription ; l'esperanto suit l'ido, qui compte juste.
% ===================================================================

%%K t10-apar-1 apar
\VUsaut{4.0mm}
\VUtitre{15.0pt}{60}{TABELO N\textsuperscript{o} 10}
\VUsaut{3.0mm}
\VUpk{11.4pt}{40}{La maro. --- La haveno.}
\VUsaut{3.0mm}

%%K t10-01-1 p
1. --- Somere, dum unuj serĉas kvieton kaj malvarmeton sur monto, aliaj
preferas iri al la marbordo. Tiel ni agis lastjare, ĉar ni tre ŝatas
\VUgras{Oceanon} \textsuperscript{(1)}.

%%K t10-02-1 p
2. --- Pri mi, mi sentas grandan plezuron kontemplante tiun imensan vastecon
de akvo kiu sin etendas ĝis la \VUgras{horizonto} \textsuperscript{(2)}, kaj
vidante foriri por longa \VUgras{transiro} la enormajn \VUgras{vaporŝipojn}
\textsuperscript{(3)} sur kiujn enŝipiĝas la vojaĝantoj irantaj en Amerikon.

%%K t10-03-1 p
3. --- Pro tio, mia patro, kiu konas mian inklinon, ĉiam elektas, por pasigi la
feriojn, iun plaĝon tre kvietan, situantan apud unu el niaj grandaj
\VUgras{militaj havenoj}.

%%K t10-03-2 p
Dum nia lasta restado, la \VUgras{eskadro} venis ankri malantaŭ la enorma
\VUgras{digo} \textsuperscript{(4)} kiu fermas kaj protektas la \VUgras{militan
havenon} \textsuperscript{(5)}, kaj mi sukcesis laŭ mia deziro admiri la
diversajn \VUgras{militajn ŝipojn}, la malgrandan \VUgras{submarŝipon}
\textsuperscript{(10)} kiu plonĝas kaj malaperas sub la akvoj kaj ankaŭ la
enorman \VUgras{kirasŝipon} \textsuperscript{(9)} kiu, ĝis nun timis preskaŭ nur
la atakojn de la \VUgras{torpedŝipo} \textsuperscript{(8)}.

%%K t10-tit-2 sub
\VUsaut{3.0mm}
\VUpk{10.2pt}{40}{La malgrandaj ŝipoj.}
\VUsaut{3.0mm}

%%K t10-04-1 p
4. --- Ĉi tiu jam sciis tre \VUgras{lerte} trovi la malfortan punkton de la dika
metalkiraso por enigi tien la danĝeran \VUgras{torpedon}, kies eksplodo kaŭzas
la pereon de la ŝipo, sed tamen ĝi estis malpli timinda ol la submarŝipo, ĉar
la kontraŭtorpedŝipoj facile persekutas ĝin.

%%K t10-05-1 p
5. --- Mi ankaŭ povis vidi la rapidajn \VUgras{krozŝipojn} \textsuperscript{(6)}
kirasitajn, preskaŭ tiel nevundeblajn kiel la vera kirasŝipo, plurajn
\VUgras{transportŝipojn} \textsuperscript{(12)} tri aŭ kvar \VUgras{kanonboatojn}
\textsuperscript{(7)} kaj fine unu \VUgras{fregaton} \textsuperscript{(11)}.
\VUgras{La spektaklo de tiu ŝiparo, kiam ĝi manovras por malankri estas plej
interesa.}

%%K t10-06-1 p
6. --- Kiom da diferenco de konstruiteco inter tiuj grandaj ŝtalamasoj kaj la
elegantaj \VUgras{trimastuloj} aŭ la graciaj \VUgras{velŝipoj}
\textsuperscript{(13)} ankoraŭ nun uzataj en la komerca ŝiparo kiun mi vidis
eniri la havenon, per ĉiuj veloj, kiam mi promenis sur la \VUgras{kajo}
\textsuperscript{(14)}. Vere ĉi tiuj perdis multon el siaj avantaĝoj
\VUgras{kiam} la \VUgras{vento} estis kontraŭa. Ili devis malhisi ĉiujn velojn
por reeniri la basenon kaj \VUgras{remorkilo} \textsuperscript{(16)} estis
necesa por venigi ilin kaj alkonduki ilin, kiel simplajn \VUgras{kargoboatojn}
\textsuperscript{(17)}, al la kajo.

%%K t10-tit-3 sub
\VUsaut{3.0mm}
\VUpk{10.2pt}{40}{La komerca haveno.}
\VUsaut{3.0mm}

%%K t10-07-1 p
7. --- Tio kio interesas en la haveno pri kiu mi parolas, estas ke ĝi enhavas ne
nur militan havenon, arte aranĝitan per gigantaj laboroj, sed ankaŭ mirindan
\VUgras{komercan havenon} situantan en la \VUgras{enfluejo} de granda fluo.

%%K t10-08-1 p
8. --- La tera lango kiu apartigas la du basenojn estas fortikigita kaj
defendata per serio da \VUgras{baterioj} kaj \VUgras{bastionoj} kiuj avancas
sur la maron. Oni bezonas specialan permeson por promeni sur la
\VUgras{remparoj} \textsuperscript{(15)}. Mi facile akiris ĝin, per mia onklo,
kiu estas inĝeniero de la \VUgras{konstruejoj} \textsuperscript{(18)} kaj mi
iris viziti la ŝipojn kiujn oni konstruas sub vastaj hangaroj.

%%K t10-09-1 p
9. --- Eĉ mi ĉeestis la lanĉon de kirasŝipo kaj mi memoras la emociigan
momenton kiun la tuta homamaso sentis kiam la grandega maso, lasita al si mem,
ekglitis sur sia lulilsimila framo kaj venis flosi sur la akvo de la fluo.

%%K t10-10-1 p
10. --- Por iri al la konstruejoj, oni trairas la \VUgras{manovrkampon}
\textsuperscript{(19)} kie la soldatoj de la garnizono ĉiutage venas ekzerci
sin, oni pasas sur fera \VUgras{ponteto} \textsuperscript{(20)}, kiu
komunikigas la paradesplanadon kun la \VUgras{arsenalo} \textsuperscript{(21)}.

%%K t10-tit-4 sub
\VUsaut{3.0mm}
\VUpk{10.2pt}{40}{La arsenalo kaj la dokoj.}
\VUsaut{3.0mm}

%%K t10-11-1 p
11. --- Kun mia onklo, mi vizitis tiun grandan konstruaĵon plenan de
\VUgras{kanonoj} \textsuperscript{(22)}, \VUgras{kugloj} \textsuperscript{(23)}
kaj \VUgras{obusoj}. Ankaŭ mi vidis la \VUgras{metalfarejon}
\textsuperscript{(24)} en kiu estas fabrikataj la \VUgras{kaldronoj}
\textsuperscript{(25)} de la vaporŝipoj.

%%K t10-12-1 p
12. --- Tre proksime, estas la \VUgras{dokoj} \textsuperscript{(26)} ---
ripardokoj --- kaj la tre rimarkindaj \VUgras{flosantaj dokoj}
\textsuperscript{(27)} en kiuj mi povis vidi tre proksime, sur riparenda ŝipo,
la \VUgras{helicon} \textsuperscript{(28)}, la \VUgras{kilon}
\textsuperscript{(29)} kaj la \VUgras{stirilon} \textsuperscript{(30)} de boato.

%%K t10-13-1 p
13. --- La komerca haveno estas tute same studinda kiel la milita haveno, kaj mi
pasigis multajn horojn gapante sur la kajoj kie estas albordigitaj la
\VUgras{radŝipoj} \textsuperscript{(31)} kiuj resekvas la fluon, kaj rigardante
funkcii la \VUgras{mastgruojn} \textsuperscript{(32)}, kiuj levas, kiel plumojn,
enormajn kargaĵojn.

%%K t10-tit-5 sub
\VUsaut{4.0mm}
\VUpk{10.2pt}{40}{La piloto.}
\VUsaut{3.0mm}

%%K t10-14-1 p
14. --- Mi admiris la \VUgras{transatlantikajn ŝipojn} \textsuperscript{(34)}
elirantajn el la rado, kun siaj \VUgras{flagoj} \textsuperscript{(35)} en plena
aero, kondukataj de lerta \VUgras{piloto} \textsuperscript{(36)} kiu mirinde
scias laŭiri la \VUgras{kanaleton} markitan per \VUgras{buoj}
\textsuperscript{(40)} kaj \VUgras{balizoj}, evitante la \VUgras{ŝarĝboatojn}
\textsuperscript{(41)} kiuj penige direktiĝas per sia unika velo kvadrata. Mi
distingis sur la \VUgras{antaŭaĵo} \textsuperscript{(37)} --- la pruo --- la
\VUgras{kabinojn} \textsuperscript{(39)} de la dua klaso, kaj sur la
\VUgras{malantaŭaĵo} \textsuperscript{(38)} --- la pobo --- la salonojn por la
pasaĝeroj de la unua klaso.

%%K t10-15-1 p
15. --- Kelkfoje mi ankaŭ ĉeestis la ŝarĝadon de \VUgras{kargoboato}
\textsuperscript{(42)} en kiun ĵus estis amasigataj la varoj de la
\VUgras{magazeno} \textsuperscript{(43)} kaj de la \VUgras{marstacio}
\textsuperscript{(44)}, alkondukitaj per la multegaj trajnoj de la
\VUgras{kajolinio} \textsuperscript{(45)}.

%%K t10-tit-6 sub
\VUsaut{3.0mm}
\VUpk{10.2pt}{40}{La remplezuro.}
\VUsaut{3.0mm}

%%K t10-16-1 p
16. --- Ofte mi boatveturis en la \VUgras{flosdoko} \textsuperscript{(53)}, en
kiun venas rifuĝi la \VUgras{barkoj} \textsuperscript{(46)}, kaj estis ĝojo por
mi manovri la \VUgras{remilon} \textsuperscript{(47)}. Mi supreniris la
\VUgras{ferdekon} \textsuperscript{(48)} de \VUgras{velbarko}
\textsuperscript{(49)}, mi grimpis, per \VUgras{ŝnuregoj}
\textsuperscript{(52)}, sur la \VUgras{maston} \textsuperscript{(51)} ĝis la
\VUgras{jardoj} \textsuperscript{(50)}, poste mi ripetis mian promenon
\VUgras{per jolo} \textsuperscript{(54)} inter pezaj \VUgras{fiŝbarkoj}
\textsuperscript{(55)}, kie la \VUgras{retoj} \textsuperscript{(56)} sekiĝas.

%%K t10-17-1 p
17. --- Sur la \VUgras{kajo} \textsuperscript{(58)} defilis antaŭ mi la plej
diversaj \VUgras{varoj} \textsuperscript{(59)}, entenataj en \VUgras{kestoj}
\textsuperscript{(60)} aŭ en \VUgras{pakegoj} \textsuperscript{(61)}. Ili formis
la \VUgras{kargaĵon} \textsuperscript{(63)} de la \VUgras{barkasoj}, kaj potencaj
\VUgras{gruoj} \textsuperscript{(62)} deprenis kaj demetis ilin sur
\VUgras{kamionojn} \textsuperscript{(64)} dum la \VUgras{kargveturistoj}
deklaris ilin kaj pagis la impostojn en la \VUgras{doganejo}
\textsuperscript{(65)}.

%%K t10-18-1 p
18. --- Fine, kiam mi alvenis al la \VUgras{turnponto} \textsuperscript{(66)} aŭ
al la \VUgras{kluzo} \textsuperscript{(69)}, pasante antaŭ la
\VUgras{dragmaŝino} \textsuperscript{(68)} kiu purigas la \VUgras{kanalon}
\textsuperscript{(67)}, mi elŝipiĝis sur la kajon apud la \VUgras{semaforo}
\textsuperscript{(70)}.

%%K t10-19-1 p
19. --- Sur la alia flanko de tiu kajo ja venas albordiĝi la \VUgras{ŝalupoj}
\textsuperscript{(71)} kiuj kondukas al la tero la oficirojn de la eskadro.
Estas admirinda spektaklo vidi la maristojn kadence levi siajn remilojn, dum
la barko portata de la \VUgras{ondo} \textsuperscript{(72)}, facile albordiĝas
al la ŝtuparo de la elŝipiĝejo. Sur ĉi tiu loko, oni pli bone povas observi la
\VUgras{tajdon} kaj la movadon de la \VUgras{fluso} kaj \VUgras{malfluso} sur
la \VUgras{plaĝo} \textsuperscript{(73)}.

%%K t10-tit-7 sub
\VUsaut{3.0mm}
\VUpk{10.2pt}{40}{La klubejo.}
\VUsaut{3.0mm}

%%K t10-20-1 p
20. --- Por reveni al la hejmo, mi uzis la \VUgras{elektran tramvojon}
\textsuperscript{(74)} kies \VUgras{haltejo} \textsuperscript{(75)} estas apud
la \VUgras{fosto} metita antaŭ la klubo de la oficiroj.

%%K t10-20-2 p
De sur la \VUgras{balkono} \textsuperscript{(76)} de la klubo, ornamita per
\VUgras{ankroj} \textsuperscript{(77)}, kanonoj kaj kugloj, ofte mi vidis
kunvenon belegan de maraj oficiroj : \VUgras{leŭtenantoj}
\textsuperscript{(78)}, kun sia spado, \VUgras{kapitanoj de artilerio}
\textsuperscript{(79)} kaj \VUgras{infanterio} \textsuperscript{(80)} kun sia
\VUgras{sabro}. Malantaŭ ili estis \VUgras{maristo} \textsuperscript{(81)} kiu
alportis al ili \VUgras{lorneton}, kiam ili volis distingi tion kio okazas
malproksime.

%%K t10-21-1 p
21. --- Mi ankaŭ vidis junajn \VUgras{subleŭtenantojn} \textsuperscript{(82)}
ricevantajn la instrukciojn de sia \VUgras{komandanto} \textsuperscript{(83)}
aŭ de la \VUgras{admiralo} \textsuperscript{(84)}, dum \VUgras{kvaronmastro}
\textsuperscript{(85)} atendis por porti ilin al la ŝipo.

%%K t10-22-1 p
22. --- De tiu balkono, la vidaĵo estas belega : oni superrigardas la tutan
plaĝon kun ĝiaj \VUgras{tendoj} \textsuperscript{(86)} kaj \VUgras{kabanoj}
\textsuperscript{(87)}, kaj oni vidas, je la horo de la banado, la
\VUgras{banantojn} \textsuperscript{(88)} kiuj ĝoje petolas antaŭ la
\VUgras{kazino} \textsuperscript{(89)}.

%%K t10-23-1 p
23. --- Ĉe la ekstremo de la plaĝo, la marbordo formas malgrandan
\VUgras{duoninsulon} \textsuperscript{(91)} \VUgras{rokan}, kie la
\VUgras{ondego} \textsuperscript{(92)} venas rompi sin kaj sur kiu estas
konstruita \VUgras{lumturo} \textsuperscript{(93)}, nokte montranta la vojon al
la navigantoj. Tiu duoninsulo komunikas kun la tero nur per \VUgras{istmo} tre
mallarĝa \textsuperscript{(90)}.

%%K t10-tit-8 sub
\VUsaut{3.0mm}
\VUpk{10.2pt}{40}{Ĝenerala vidaĵo de la marbordoj.}
\VUsaut{3.0mm}

%%K t10-24-1 p
24. --- La \VUgras{klifo} \textsuperscript{(94)} etendas sin, disfendita de la
maro, kiu kaprice desegnas en ĝi serion da \VUgras{golfetoj}
\textsuperscript{(95)} kaj \VUgras{promontoroj} (kaboj), igante ĝin plej
pentrinda.

%%K t10-24-2 p
Sur la \VUgras{altebenaĵo} \textsuperscript{(96)}, kiu superregas la
\VUgras{(mar)markolon} \textsuperscript{(98)}, estas \VUgras{forto}
\textsuperscript{(97)} tre grava kaj kelkaj « vilaoj ».

%%K t10-25-1 p
25. --- Tiu markolo apartigas de la kontinento \VUgras{insulon}
\textsuperscript{(99)} tre danĝeran, ĉirkaŭitan de \VUgras{rifoj}
\textsuperscript{(100)} kaj \VUgras{ondorompaĵoj}, kie barkoj kaj eĉ grandaj
ŝipoj surterigas iufoje. Malgraŭ la pretemo de la helpoj, kaj la rapida alveno
de la \VUgras{savboatoj}, tiuj \VUgras{ŝippereoj} estas terurigaj, kaj, dum la
\VUgras{tempestoj} ekvinoksaj, pluraj ŝipoj pereis kun homoj kaj varoj.

%%K t10-25-2 p
Malgraŭ tiu najbaraĵo, \VUgras{la haveno} estas tre vizitata de la maristoj.
\VUsaut{6.0mm}
\VUfilet{22mm}
